\akhead{Topic 9 — Exponents, Radicals, and Rational Exponents}

\ans{1}{1}{A}{Same base, so add the exponents: $x^{4}\cdot x^{7}=x^{4+7}=x^{11}$.}

\ans{1}{2}{$9$}{Any nonzero base to the power $0$ is $1$, so $7^{0}+2^{3}=1+8=9$.}

\ans{1}{3}{A}{Same base, so subtract the exponents:
$\dfrac{x^{9}}{x^{4}}=x^{9-4}=x^{5}$.}

\ans{1}{4}{$\dfrac{1}{16}$}{A negative exponent means a reciprocal:
$4^{-2}=\dfrac{1}{4^{2}}=\dfrac{1}{16}$.}

\ans{1}{5}{A}{A power raised to a power multiplies the exponents:
$\left(x^{3}\right)^{5}=x^{3\cdot 5}=x^{15}$.}

\ans{1}{6}{$11$}{$\sqrt{49}=7$ and $\sqrt{16}=4$, so the sum is $7+4=11$.}

\ans{1}{7}{A}{An $n$th root is the exponent $\tfrac{1}{n}$, so
$\sqrt[3]{x}=x^{1/3}$.}

\ans{1}{8}{A}{Multiply the leading factors and add the powers of ten:
$(3\cdot 2)\times 10^{4+3}=6\times 10^{7}$.}

\ans{1}{9}{$4$}{$\sqrt{48}=\sqrt{16\cdot 3}=\sqrt{16}\,\sqrt{3}=4\sqrt{3}$, so
$a=4$.}

\ans{1}{10}{A}{The exponent applies to both factors:
$\left(2x^{3}\right)^{4}=2^{4}x^{3\cdot 4}=16x^{12}$.}

\ans{1}{11}{A}{Pull out the perfect square:
$\sqrt{50}=\sqrt{25\cdot 2}=5\sqrt{2}$.}

\ans{1}{12}{A}{Write $32=2^{5}$. Then $2^{x+1}=2^{5}$, so $x+1=5$ and $x=4$.}

\ans{1}{13}{$27$}{$9^{3/2}=\left(\sqrt{9}\right)^{3}=3^{3}=27$.}

\ans{1}{14}{A}{Multiply above and below by $\sqrt{3}$:
$\dfrac{6}{\sqrt{3}}\cdot\dfrac{\sqrt{3}}{\sqrt{3}}=\dfrac{6\sqrt{3}}{3}=2\sqrt{3}$.}

\ans{1}{15}{A}{Each factor crossing the bar flips the sign of its exponent:
$\dfrac{3x^{-2}}{y^{-3}}=\dfrac{3y^{3}}{x^{2}}$.}

\ans{1}{16}{A}{$\left(2.5\times 10^{-5}\right)\left(4\times 10^{6}\right)
=10\times 10^{1}=100$ grams.}

\ans{1}{17}{$9$}{$64^{2/3}=\left(\sqrt[3]{64}\right)^{2}=4^{2}=16$ and
$49^{1/2}=7$, so the value is $16-7=9$.}

\ans{1}{18}{A}{Write $27=3^{3}$: $3^{2x}=3^{3(x-1)}=3^{3x-3}$. Equate
exponents: $2x=3x-3$, so $x=3$.}

\ans{1}{19}{A}{In exponent form the quotient is
$\dfrac{x^{5/3}}{x^{1/2}}=x^{5/3-1/2}=x^{7/6}$.}

\ans{1}{20}{A}{Multiply by the conjugate $\sqrt{5}+1$:
$\dfrac{4\left(\sqrt{5}+1\right)}{5-1}=\dfrac{4\left(\sqrt{5}+1\right)}{4}
=\sqrt{5}+1$.}

\ans{1}{21}{$\dfrac{15}{4}$}{Over base $2$: $2^{-x}=2^{3(x-5)}$, so
$-x=3x-15$, giving $4x=15$ and $x=\tfrac{15}{4}$.}

\ans{1}{22}{A}{Subtract exponents on each base: $a^{-3-2}b^{5-(-1)}=a^{-5}b^{6}
=\dfrac{b^{6}}{a^{5}}$.}

\ans{2}{1}{A}{Cube every factor:
$\left(3x^{2}y^{3}\right)^{3}=3^{3}x^{2\cdot 3}y^{3\cdot 3}=27x^{6}y^{9}$.}

\ans{2}{2}{A}{$3\sqrt{8}=3\cdot 2\sqrt{2}=6\sqrt{2}$ and $\sqrt{50}=5\sqrt{2}$,
so the difference is $6\sqrt{2}-5\sqrt{2}=\sqrt{2}$.}

\ans{2}{3}{$\dfrac{3}{2}$}{Write both sides over base $2$: $2^{2x}=2^{3}$, so
$2x=3$ and $x=\tfrac{3}{2}$.}

\ans{2}{4}{A}{Divide the leading factors and subtract the powers:
$\dfrac{8.4}{2}\times 10^{9-4}=4.2\times 10^{5}$.}

\ans{2}{5}{A}{A negative outer exponent inverts the fraction first:
$\left(\dfrac{y^{2}}{x^{3}}\right)^{2}=\dfrac{y^{4}}{x^{6}}$.}

\ans{2}{6}{$3$}{$\sqrt{18x^{4}}=\sqrt{9x^{4}}\,\sqrt{2}=3x^{2}\sqrt{2}$ for
$x>0$, so $a=3$.}

\ans{2}{7}{A}{The side is $\sqrt{128}=\sqrt{64\cdot 2}=8\sqrt{2}$ centimetres.}

\ans{2}{8}{A}{Add the fractional exponents:
$x^{1/2+1/3}=x^{3/6+2/6}=x^{5/6}$.}

\ans{2}{9}{$\dfrac{9}{2}$}{Write $9=3^{2}$: $3^{2(x-1)}=3^{7}$, so $2x-2=7$ and
$x=\tfrac{9}{2}$.}

\ans{2}{10}{A}{$\left(2\times 10^{9}\right)\left(3\times 10^{-3}\right)
=6\times 10^{9-3}=6\times 10^{6}$ operations.}

\ans{2}{11}{A}{$\sqrt{20}=2\sqrt{5}$, so the expression is
$\dfrac{10}{2\sqrt{5}}=\dfrac{5}{\sqrt{5}}=\sqrt{5}$.}

\ans{2}{12}{$\dfrac{4}{9}$}{The negative exponent inverts the fraction:
$\left(\dfrac{8}{27}\right)^{2/3}=\left(\dfrac{2}{3}\right)^{2}=\dfrac{4}{9}$.}

\ans{2}{13}{A}{Multiply each exponent by $4$:
$x^{(1/2)\cdot 4}y^{(-1/4)\cdot 4}=x^{2}y^{-1}=\dfrac{x^{2}}{y}$.}

\ans{2}{14}{A}{$f(64)=64^{2/3}=\left(\sqrt[3]{64}\right)^{2}=4^{2}=16$.}

\ans{2}{15}{A}{The product gives $a+b=12$. Substituting $a=2b$ gives $3b=12$,
so $b=4$.}

\ans{2}{16}{$\dfrac{13}{3}$}{Over base $2$:
$2^{x}\cdot 2^{2(x+1)}=2^{3x+2}$ and $8^{5}=2^{15}$, so $3x+2=15$ and
$x=\tfrac{13}{3}$.}

\ans{2}{17}{A}{Multiply above and below by the conjugate $\sqrt{x}+3$:
$\dfrac{\left(\sqrt{x}+3\right)^{2}}{x-9}=\dfrac{x+6\sqrt{x}+9}{x-9}$.}

\ans{2}{18}{$5$}{The left side is $x^{3k}\cdot x^{2}=x^{3k+2}$, so $3k+2=17$
and $k=5$.}

\ans{2}{19}{A}{Since $2^{n}3^{n}=6^{n}$, the denominator is
$\dfrac{6^{n}}{3}$, so the quotient is
$\dfrac{6\cdot 6^{n}\cdot 3}{6^{n}}=18$.}

\ans{2}{20}{A}{The numerator is $16\times 10^{-6}$, so the quotient is
$\dfrac{16}{8}\times 10^{-6-(-5)}=2\times 10^{-1}$.}

\ans{2}{21}{A}{$\sqrt{c}\cdot c^{3/2}=c^{1/2+3/2}=c^{2}$, so $c^{2}=64$ and,
since $c>0$, $c=8$.}

\ans{2}{22}{$-2$}{Inside the parentheses $a^{2-(-1)}b^{-1-3}=a^{3}b^{-4}$;
squaring gives $a^{6}b^{-8}$, so $m+n=6+(-8)=-2$.}
